\documentclass{article}
\usepackage[utf8]{inputenc}
\usepackage{amsmath}
\usepackage{geometry}
\geometry{margin=2cm}
\begin{document}

\section*{Questão 106 — ENEM 2024 — Ciências da Natureza}

\subsection*{Interpretação da Questão}
A questão apresenta uma situação em que Calvin, depois de se soltar de um balanço, cai ao chão. No momento do impacto, o chão exerce uma força sobre ele, chamada força normal. Deve-se identificar qual vetor entre as alternativas melhor representa essa força resultante.

\subsection*{Estratégia para Resolução}
Para resolver a questão, devemos analisar a trajetória de Calvin:
\begin{itemize}
  \item Inicialmente, Calvin está em movimento circular preso ao balanço.
  \item Ao se soltar, seu movimento passa a ser parabólico, sob ação da gravidade.
  \item Durante a queda, a única força que atua é a gravidade, para baixo.
  \item No instante de contato com o solo, o chão exerce uma força normal para cima (reação ao peso), evitando que Calvin penetre a superfície.
  \item Considerando que Calvin conserva componente de velocidade horizontal no impacto, o chão exerce, além da força vertical, uma componente horizontal contrária ao movimento para desacelerar Calvin.
  \item Comparar os vetores das alternativas para escolher o que corresponde a essa força resultante.
\end{itemize}

\subsection*{Análise dos Conteúdos Envolvidos}
\begin{center}
\begin{tabular}{|p{4cm}|p{5cm}|p{6cm}|}
\hline
\textbf{Conceito} & \textbf{Ideia Principal} & \textbf{Aplicação à Questão} \\
\hline
Trajetória parabólica & Movimento sob efeito da gravidade, formando uma curva parabólica & Calvin descreve essa trajetória após soltar o balanço \\
\hline
Força normal & Força de reação do solo, perpendicular à superfície & No impacto, exerce força para cima e contra o movimento horizontal \\
\hline
Composição vetorial das forças & Forças decompostas em componentes vertical e horizontal & Identificar vetor resultante da soma dessas forças no impacto \\
\hline
\end{tabular}
\end{center}

\subsection*{Teoria Específica}
\begin{itemize}
  \item Corpos soltos em um ponto e submetidos à gravidade seguem uma trajetória parabólica, com movimento horizontal constante e movimento vertical uniformemente acelerado.
  \item A força normal é uma força de reação perpendicular à superfície de apoio, que impede penetração.
  \item Essa força pode ter componentes além da vertical para modificar o movimento horizontal do corpo.
  \item A decomposição vetorial é necessária para interpretar corretamente direção e sentido da força.
\end{itemize}

\subsection*{Resolução Detalhada}
\begin{itemize}
  \item Calvin solta o balanço: passa de movimento circular para trajeto parabólico.
  \item Durante a queda, atua apenas a gravidade ($\vec{F}_g$) para baixo.
  \item No impacto, ocorre a interação entre Calvin e o solo, que exerce força normal ($\vec{F}_n$).
  \item A força normal tem componente vertical para cima para sustentar o peso.
  \item Como Calvin ainda possui velocidade horizontal para a direita, o solo exerce também componente horizontal para a direita, contrária ao movimento, para desacelerar.
  \item Portanto, a força resultante da interação é um vetor diagonal para cima e para a direita.
  \item Dentre as alternativas, apenas o vetor da alternativa \textbf{A} representa essa força combinada.
\end{itemize}

\subsection*{Pulo do Gato}
No instante do impacto, a força do chão sempre apresenta componente vertical para cima (força normal) e pode ter uma componente horizontal contrária à velocidade para desaceleração. A alternativa que melhor combina essas duas é correta.

\subsection*{Armadilhas Comuns}
\begin{itemize}
  \item\, Confundir direção da força resultante com a trajetória do movimento.
  \item\, Ignorar a componente horizontal da força normal.
  \item\, Pensar que a força resultante é somente vertical ou somente horizontal.
  \item\, Confundir força normal com força gravitacional ou atrito.
\end{itemize}

\subsection*{Código da Questão}
\texttt{CIENCIAS_DA_NATUREZA_FISICA_001}

\vspace{1cm}
\centerline{\rule{12cm}{0.4pt}}

\begin{flushright}
\emph{Resolução aprovada e consistente.}
\end{flushright}


\end{document}